\documentclass[11pt,a4paper]{article}
\usepackage[utf8]{inputenc}
\usepackage[T1]{fontenc}
\usepackage[spanish]{babel}
\usepackage{hyperref}
\usepackage{geometry}
\geometry{margin=2.5cm}
\usepackage{parskip}

\title{114 un autor tiene N candados, y cada uno se abre por separado}
\author{Alejandro Valenzuela Rivera\\\small Universidad Autónoma de Nuevo León}
\date{\today}
\begin{document}
\maketitle

\begin{quote}
\textbf{Sesgo declarado.} Este trabajo se sostiene desde el sesgo declarado de su autor, y no finge neutralidad. Lo que afirma se puede tumbar, y ese es el punto: el falsificador está escrito, no escondido.
\end{quote}

\section*{Resumen}
Cada posicion que el autor sostuvo y que puede ser atacada por su cuenta es un candado: un archivo corto con lo que sostiene, por que, de donde viene, y como se abre. La tesis entera es el candado cero, del que salen los demas.

\section{Por qué se sostiene}
Porque una voz con un solo candado es una voz que se calla lo demas que piensa. Y porque una tesis que solo se puede atacar como un bloque no se puede contestar: se acepta o se rechaza entera. Un candado corto se abre con una llave y deja el resto de pie.

\section{De dónde viene}
La autoridad, 2026-10-26.

\section{El caso}
Ninguna: es la regla del candado.

\section{Método}
N candados que se abren uno por uno, cada uno con su friccion.

\section{Límites, y qué lo tumba}
Un candado que se abre arrastrando a otros, o uno que no se abre nunca.

\section{Lo que queda abierto}
\textquestiondown{}Cuantos cortes tiene una idea antes de que deje de ser una?

\section*{Cómo se leyó esta tesis}
Esto se escribió en el repositorio man-made-empathy, del que este candado es una pieza: candados/01 el ente y el metodo/114 un autor tiene N candados, y cada uno se abre por separado.md. La máquina que lo compuso lo hizo desde este mismo archivo, y no agregó nada que no estuviera escrito. Si algo de esto es falso, se dice en el mismo sitio, y se corrige con fecha.

\end{document}
